\section{The newly certified native figure-eight cut}
\label{sec:native-target}

The retained fixture \code{pilot/finite\_figureEight/0} provides a concrete
source-authenticated example that crosses the preceding cancellation barrier.
It contains ten tetrahedra and 36 normal disks. Its finite core has
\[
 (h_0,h_1,h_2,h_3)=(45,68,17,0).
\]
There are five ported parallelity bases, all disks, and every one of their
vertical boundaries is attached as a full annulus. The source cut Euler
characteristic is $-1$. The whole-cut compiler produces two presentations,
with $(a,m,L)=(7,5,18)$ and $(18,17,35)$, respectively. Thus the aggregate
presentation has 25 generators, 22 relators and 53 letters.

\subsection{A complete hand-checkable reduction}

For the difficult component, rename the native generators as
\[
 a=33,\ b=41,\ c=45,\ d=49,\ e=57,\ f=61,\ g=65.
\]
In native row order 17 through 21, the five annular relations are
\begin{align*}
 &gcf^{-1}b^{-1}f^{-1}d^{-1}, & &bed^{-1},\\
 &e^{-1}a, & &adgc^{-1},\\
 &egf^{-1}. &&
\end{align*}
They are already freely and cyclically reduced. Every generator has at least
two occurrences across the family, so there is no initial unsigned column
leaf. Therefore merely normalizing words before invoking the old peeling
routine does not resolve this component.

The following sequence gives a transparent complete algebraic proof:
\begin{enumerate}
\item From $e^{-1}a=1$, eliminate $a$ by $a=e$.
\item From $bed^{-1}=1$, eliminate $b$ by $b=de^{-1}$.
\item From $egf^{-1}=1$, eliminate $e$ by $e=fg^{-1}$.
\item The fourth relation becomes $fg^{-1}dgc^{-1}=1$; eliminate
      $c=fg^{-1}dg$.
\item The first relation then freely reduces to
      $gfg^{-1}f^{-1}d^{-1}=1$; eliminate $d=gfg^{-1}f^{-1}$.
\end{enumerate}
Only $f,g$ survive and no relation remains. Each step solves a generator that
occurs once in its donor. The total reduced lengths are
\[
 18\longrightarrow16\longrightarrow14\longrightarrow11
 \longrightarrow5\longrightarrow0.
\]
No freely trivial relator is discarded separately. This proves the first
component group is $F_2$ without trusting a group-recognition heuristic.

The deterministic production policy chooses a shorter intermediate trace:
\[
 (17,41),\ (21,57),\ (18,45),\ (20,49),\ (19,33),
\]
where a pair records the donor row and eliminated generator. Its length
sequence is
\[
 18\longrightarrow16\longrightarrow13\longrightarrow7
 \longrightarrow3\longrightarrow0.
\]
It uses three substitutions involving other relators and then two private
column pivots. The independent checker reconstructs every substitution and
cyclic reduction from the original source-bound words.

The second component needs seventeen private pivots and leaves one generator,
native label 68. Its free rank is one. The aggregate values are consequently
\[
 c=2,\qquad r=2+1=3,\qquad \chi(W)=-1=c-r.
\]
Theorem~\ref{thm:free-euler} certifies a genus-two handlebody and a solid torus.
Here $h_3=z=0$, so the relator ledger agrees automatically.

\subsection{What this experiment establishes}

The original normal coordinates, source triangulation, complete compressed
proof packet and exact row data are included. The public terminal verifier
and the dedicated new verifier both replay the new schema. A separately
expanded Regina cut has 319 tetrahedra, split into components of 191 and 128
tetrahedra; its recognition gives genera one and two. This comparison supplies
an external implementation check, while the delivered certificate verifier
uses the finite source charts and theorems above, without invoking Regina.

The incoming deferred audit identifies this fixture as the remaining actual
handlebody cut missed by that method~\cite{DeferredForests}. The new result is
a strict increase in native whole-cut coverage. It is not an example of
recognizing the figure-eight knot as an unknot. It also does not establish that
all handlebody cuts have nonincreasing singleton traces: the explicit
automorphism counterexample in Section~\ref{sec:no-growth-tietze} prevents
that inference.
